%% ============================================================
%%  Chapitre 1 — Introduction
%% ============================================================
\chapter{Introduction}
\label{ch:introduction}

\section{Contexte et motivation}

La production scientifique mondiale connaît une croissance exponentielle :
selon les estimations, plus de deux millions d'articles de recherche sont
publiés chaque année dans les seules disciplines relevant des sciences et
technologies de l'information \cite{priem2022openalex}. Pour un chercheur
débutant sa thèse ou souhaitant explorer un domaine adjacent au sien, cette
abondance constitue paradoxalement un obstacle. Il est désormais impossible
de lire l'intégralité de la littérature pertinente ; la sélection, la
hiérarchisation et la synthèse des travaux existants sont devenues des tâches
à part entière, chronophages et susceptibles d'erreurs d'omission.

Les pratiques actuelles de veille scientifique reposent principalement sur des
approches manuelles ou semi-automatisées : interrogation des moteurs de
recherche académiques (Google Scholar, Semantic Scholar, Web of Science),
lecture d'articles de revue (\textit{survey papers}), alertes par mots-clés,
ou recommandation par les pairs. Ces approches présentent plusieurs limitations
majeures. Premièrement, elles requièrent un investissement temporel considérable,
estimé à plusieurs semaines pour une revue de littérature exhaustive dans un
nouveau domaine \cite{kitchenham2007guidelines}. Deuxièmement, elles sont
sujettes à des biais de sélection, en particulier pour les chercheurs non
initiés qui manquent des critères experts pour évaluer la qualité d'un travail.
Troisièmement, elles ne permettent pas d'actualisation automatique lorsque de
nouveaux articles paraissent.

L'émergence des grands modèles de langage (LLM, \textit{Large Language Models})
\cite{openai2023gpt4, anthropic2024claude} offre de nouvelles perspectives pour
l'automatisation de ces tâches cognitives. La capacité des LLM à comprendre,
résumer et synthétiser des textes techniques ouvre la voie à des systèmes
capables d'assister le chercheur de manière substantielle. Parallèlement,
les architectures multi-agents fondées sur des LLM \cite{wang2024survey}
permettent de décomposer des tâches complexes en sous-tâches spécialisées,
d'implémenter des mécanismes de contrôle qualité, et d'exploiter les forces
complémentaires de plusieurs modèles ou stratégies de raisonnement.

\section{Problématique}

La problématique centrale de ce mémoire peut être formulée en deux questions
imbriquées :

\begin{enumerate}
  \item \textbf{Comment concevoir un système automatisé de veille scientifique}
    qui produise des résultats de qualité comparable à celle d'un expert humain,
    en termes de pertinence des articles sélectionnés, de fidélité des résumés
    et de richesse de l'analyse des tendances ?
  \item \textbf{Comment rendre ce système généralisable à de nouveaux domaines}
    sans nécessiter une reconfiguration manuelle coûteuse pour chaque nouveau
    domaine de recherche ?
\end{enumerate}

La seconde question est au cœur de la contribution scientifique de ce travail.
Un système de veille conçu spécifiquement pour un domaine (par exemple, la
détection de fausses nouvelles) exploitera des critères de qualité spécifiques
à ce domaine : l'importance de la vélocité des citations pour les sujets
d'actualité, ou le poids accordé à la pertinence sémantique pour des champs
très techniques. Transférer naïvement ces critères à un autre domaine (par
exemple, la biologie computationnelle) peut conduire à des résultats
sous-optimaux. La généralisation constitue donc un enjeu technique et
scientifique de premier ordre.

\section{Objectifs}

Ce mémoire poursuit les objectifs suivants :

\begin{enumerate}
  \item \textbf{Concevoir et implémenter} un pipeline multi-agents complet de
    veille scientifique, intégrant la recherche d'articles, le filtrage par
    qualité, la summarisation, la synthèse et l'analyse des tendances.

  \item \textbf{Optimiser automatiquement} les paramètres de scoring documentaire
    par domaine, en utilisant des techniques d'AutoML (Optuna TPE), et démontrer
    l'amélioration par rapport aux configurations expertes fixes.

  \item \textbf{Développer une approche de méta-learning} qui prédit les
    paramètres optimaux pour un nouveau domaine à partir de caractéristiques
    mesurables de ce domaine (\textit{meta-features}), sans nécessiter
    d'optimisation préalable dans ce domaine.

  \item \textbf{Évaluer rigoureusement} la généralisation cross-domaine sur
    19 domaines de recherche couvrant à la fois les domaines de l'intelligence
    artificielle et des disciplines hors-IA (biologie computationnelle,
    traitement du langage médical, finance quantitative).
\end{enumerate}

\section{Hypothèses de recherche}
\label{sec:hypotheses}

Cinq hypothèses structurent l'évaluation expérimentale de ce mémoire.
Parmi celles-ci, deux ont été confirmées et constituent les résultats
centraux du travail :

\begin{hypothese}[H1 — AutoML par domaine]
  Les pondérations apprises par Optuna sur un domaine donné surpassent
  les pondérations expertes fixes sur ce même domaine, avec une amélioration
  relative d'au moins 5\,\% en NDCG@15.
\end{hypothese}

\begin{hypothese}[H3 — Méta-learning cross-domaine]
  La prédiction des pondérations optimales pour un nouveau domaine via
  une fonction $f_\theta : \text{meta\_features}(D) \to \text{weights}$,
  apprise sur d'autres domaines, surpasse le Transfer Direct
  (utilisation des pondérations moyennes) sur au moins 75\,\% des
  domaines de test.
\end{hypothese}

Les hypothèses H2, H4 et H5 portent respectivement sur la dégradation
prévisible du Transfer Direct, l'identification des goulots de généralisation
par l'architecture multi-agents, et la comparaison des profils de généralisation
de Claude versus GPT. Ces dernières hypothèses feront l'objet d'investigations
dans les travaux futurs (Chapitre~\ref{ch:conclusion}).

\section{Contributions}

Les contributions de ce mémoire sont les suivantes :

\paragraph{Contribution 1 — Système de veille scientifique opérationnel.}
Nous présentons un système complet, généraliste et déployable, fondé sur
sept agents LLM orchestrés par LangGraph, intégrant quatre fournisseurs
de LLM (Anthropic, OpenAI, Groq, Nebius) via une abstraction unifiée.

\paragraph{Contribution 2 — Scorer de qualité multi-dimensionnel et AutoML.}
Nous proposons un scorer à six dimensions calibré par domaine via
l'algorithme TPE d'Optuna, atteignant une amélioration moyenne de
$+68.7\,\%$ en NDCG@15 par rapport aux pondérations par défaut (H1 confirmée).

\paragraph{Contribution 3 — Méta-learning pour la généralisation cross-domaine.}
Nous montrons qu'un ensemble de méta-learners (\textit{Borda Count} combinant
kNN et BayesianRidge) prédit efficacement les pondérations optimales pour
de nouveaux domaines, surpassant le Transfer Direct sur 15 domaines sur 19
(79\,\%, test de Wilcoxon : $W^+ = 149$, $p = 0.014$, effet large $r_{rb} = +0.57$).

\paragraph{Contribution 4 — Patterns d'agents avancés.}
Nous appliquons les patterns Reflexion \cite{shinn2023reflexion} et
ReAct \cite{yao2022react} au contexte de la veille scientifique, et
évaluons l'optimisation automatique des prompts via ProTeGi
\cite{pryzant2023automatic}.

\section{Organisation du mémoire}

Ce mémoire est organisé comme suit :

\begin{itemize}
  \item Le \textbf{Chapitre~\ref{ch:etat_art}} présente l'état de l'art en
    matière de systèmes multi-agents LLM, d'AutoML pour la configuration
    de systèmes, de méta-learning, et de veille scientifique automatisée.

  \item Le \textbf{Chapitre~\ref{ch:methodologie}} décrit la méthodologie
    générale : la source de données (OpenAlex), les six dimensions de qualité,
    les oracles d'évaluation et le protocole expérimental.

  \item Le \textbf{Chapitre~\ref{ch:automl}} présente le scorer de qualité
    et son optimisation par AutoML, avec les résultats de l'hypothèse H1
    sur 19 domaines.

  \item Le \textbf{Chapitre~\ref{ch:agents}} décrit l'architecture
    multi-agents : pipeline LangGraph, contrat de communication Pydantic,
    et les sept agents spécialisés.

  \item Le \textbf{Chapitre~\ref{ch:patterns}} détaille les patterns avancés :
    ReAct (expansion de requêtes), Reflexion (révision de synthèse), et
    ProTeGi (optimisation de prompts).

  \item Le \textbf{Chapitre~\ref{ch:metalearning}} présente la contribution
    centrale : l'extraction de meta-features, les méta-learners évalués,
    et la validation de l'hypothèse H3 par test de Wilcoxon.

  \item Le dernier chapitre conclut et propose des pistes pour les
    travaux futurs, notamment la Phase 7 (agents concurrents Best-of-N)
    et le renforcement du protocole expérimental.
\end{itemize}
